% Pattern: proof/wrapper. A proof wrapper whose first argument is optional, as
% \newenvironment{Proof}[1][\proofname]{\begin{proof}[#1]}, takes the option of
% the use and no letter of the body; a wrapper whose second argument names the
% statement in the end mark, as \begin{PROOF}{\ref{key}}, proves it; and a
% label that a command of the paper sets, as \prlabel{key}, is a label of the
% statement.
% seen in the gold pilot (set X)
% want thm-1 statement/theorem 1 labels=th:a
% want prop-1 statement/proposition 2 labels=pr:b
% want lem-1 statement/lemma 3 labels=c
% want proof-1 proof - proves=thm-1
% want proof-2 proof - proves=lem-1
% want proof-3 proof - proves=prop-1
\documentclass{amsart}
\newtheorem{thm}{Theorem}
\newtheorem{prop}[thm]{Proposition}
\newtheorem{lem}[thm]{Lemma}
\newcommand{\prlabel}[1]{\label{pr:#1}}
\newcommand{\prref}[1]{Proposition~\ref{pr:#1}}
\newenvironment{Proof}[1][\proofname]{\begin{proof}[#1]}{\end{proof}}
\newenvironment{PROOF}[2][\proofname.]
   {\begin{proof}[#1]\renewcommand{\qedsymbol}{$\square_{#2}$}}
   {\end{proof}}
\begin{document}
\begin{thm}\label{th:a} Each tree with $n$ vertices has $n-1$ edges. \end{thm}
\begin{prop}\prlabel{b} Each forest is a union of trees. \end{prop}
\begin{lem}\label{c} Each tree with two vertices or more has a leaf. \end{lem}

\begin{Proof}[Proof of Theorem~\ref{th:a}]
Let a tree be given. \end{Proof}

\begin{PROOF}{\ref{c}} Take a longest path. \end{PROOF}

\renewcommand{\proofname}{Proof of \prref{b}}
\begin{proof} Take the components. \end{proof}
\end{document}
